%--------------------------------------------------------------------%
%
% Kata Pengantar (Pedoman UM Lampiran 7)
%
% DRAF: isi bagian dalam kurung siku dengan data sebenarnya, lalu hapus
% tanda kurung siku. Urutan ucapan terima kasih mengikuti pedoman, yaitu
% pembimbing dan program studi lebih dahulu, keluarga kemudian.
%
%--------------------------------------------------------------------%

\chapter*{Kata Pengantar}
\phantomsection
\addcontentsline{toc}{chapter}{KATA PENGANTAR}

Puji syukur penulis panjatkan kepada Tuhan Yang Maha Esa atas rahmat dan karunia-Nya. Atas kehendak-Nya penulis dapat menyelesaikan skripsi yang berjudul ``\judul''. Skripsi ini disusun untuk memenuhi salah satu persyaratan dalam menyelesaikan program \gelarprogram.

Penyusunan skripsi ini berlangsung melalui proses panjang. Pada setiap tahapnya penulis memperoleh bantuan, arahan, dan dukungan dari berbagai pihak. Karena itu, pada kesempatan ini penulis menyampaikan terima kasih kepada pihak-pihak yang berperan dalam penyelesaian skripsi ini.

\begin{enumerate}
    \item \pembimbingSatu, selaku dosen pembimbing, atas arahan, ketelitian, dan waktu yang diberikan. Bimbingan beliau membantu penulis merumuskan masalah penelitian, menyusun metode, dan menafsirkan hasil pengujian.
    \item \dekan, selaku Dekan \fakultas, atas kebijakan dan fasilitas yang mendukung pelaksanaan penelitian ini.
    \item \ketuaDepartemen, selaku Ketua Departemen \departemen, atas dukungan administratif selama masa studi penulis.
    \item Bapak Dr. Eng. Ir. Didik Dwi Prasetya, S.T., M.T., selaku Ketua \programstudi\ Universitas Negeri Malang, atas arahan mengenai ketentuan penulisan skripsi.
    \item \pembimbingSatu, selaku Dosen Pembimbing Akademik, atas nasihat dan motivasi selama masa studi.
\end{enumerate}

Penulis menyadari bahwa skripsi ini masih memiliki kekurangan. Karena itu, kritik dan saran yang membangun penulis harapkan untuk perbaikan pada penelitian selanjutnya. Penulis berharap skripsi ini bermanfaat bagi pengembangan pemrosesan bahasa alami berbahasa Indonesia.

\vspace{2\baselineskip}
\begin{flushright}
\begin{tabular}{c}
\kotatanggal \\
Penulis
\end{tabular}
\end{flushright}
